% =====================================================================
%  PARTIE 2 — Principes techniques du TinyML
% =====================================================================
\chapter{Principes techniques du TinyML}
\label{ch:principes}

\section{Introduction}

La partie précédente a permis de définir le concept de TinyML et d'en
identifier les enjeux fondamentaux. Cette partie examine les mécanismes
techniques sous-jacents rendant possible l'exécution de modèles d'intelligence
artificielle sur des dispositifs embarqués aux ressources limitées.

La problématique centrale s'énonce ainsi : comment déployer un modèle
d'apprentissage automatique, dont la taille peut atteindre plusieurs centaines
de mégaoctets, sur un microcontrôleur disposant seulement de quelques centaines
de kilooctets de mémoire vive ? La réponse repose sur trois piliers : un
pipeline de développement rigoureusement structuré, l'application de techniques
avancées de compression des modèles, et une exécution optimisée sur le matériel
cible.

% ---------------------------------------------------------------------
\section{Pipeline de développement TinyML}

Le TinyML ne consiste pas en l'entraînement direct des modèles sur le
microcontrôleur. Le dispositif embarqué n'effectue qu'une seule opération :
l'\textbf{inférence}, c'est-à-dire l'exécution d'un modèle préalablement
entraîné et optimisé. L'intégralité des calculs intensifs s'effectue en amont,
sur des stations de travail ou dans le cloud~\cite{warden2019tinyml}.

\begin{figure}[H]
\centering
\begin{tikzpicture}[node distance=0.45cm]
  \node[bloc, minimum width=2.1cm] (a) {Collecte\\\footnotesize(données)};
  \node[bloc, minimum width=2.1cm, right=of a] (b) {Prétraitement\\\footnotesize(nettoyage)};
  \node[bloco, minimum width=2.1cm, right=of b] (c) {Entraînement\\\footnotesize(PC / cloud)};
  \node[bloco, minimum width=2.1cm, right=of c] (d) {Optimisation\\\footnotesize(compression)};
  \node[blocv, minimum width=2.1cm, below=1.1cm of d] (e) {Déploiement\\\footnotesize(\code{.tflite} $\to$ C)};
  \node[blocv, minimum width=2.1cm, left=of e] (f) {Inférence\\\footnotesize(MCU)};
  \node[blocv, minimum width=2.1cm, left=of f] (g) {Action\\\footnotesize(résultat)};
  \draw[fleche] (a) -- (b);
  \draw[fleche] (b) -- (c);
  \draw[fleche] (c) -- (d);
  \draw[fleche] (d) -- (e);
  \draw[fleche] (e) -- (f);
  \draw[fleche] (f) -- (g);
  \begin{scope}[on background layer]
    \node[draw=orange, dashed, rounded corners=6pt, fit=(a)(d), inner sep=6pt,
          label={[font=\footnotesize\itshape,text=orange]above:sur ordinateur ou dans le cloud}] {};
    \node[draw=vert, dashed, rounded corners=6pt, fit=(g)(e), inner sep=6pt,
          label={[font=\footnotesize\itshape,text=vert]below:sur le microcontrôleur}] {};
  \end{scope}
\end{tikzpicture}
\caption{Pipeline complet de développement TinyML.}
\label{fig:pipeline}
\end{figure}

Le processus se décompose en cinq étapes séquentielles :
\begin{enumerate}
  \item \textbf{Collecte et préparation des données} : l'acquisition de données
        représentatives constitue l'étape fondatrice. Les données proviennent
        généralement des capteurs du dispositif final. Cette phase représente
        une part majeure du temps total de développement, car la qualité d'un
        modèle est directement dépendante de la qualité des données utilisées.
  \item \textbf{Entraînement du modèle} : cette phase s'effectue sur des
        machines dotées de ressources computationnelles importantes, utilisant
        des frameworks tels que TensorFlow ou PyTorch. Le modèle produit est
        généralement trop volumineux (plusieurs mégaoctets en Float32) pour les
        contraintes des microcontrôleurs.
  \item \textbf{Optimisation du modèle} : étape spécifique et critique du
        TinyML, elle applique des techniques de compression (quantification,
        élagage, distillation) pour réduire la taille tout en préservant les
        performances. Le modèle est converti dans le format FlatBuffer de
        TensorFlow Lite (\code{.tflite}).
  \item \textbf{Intégration dans le firmware} : le modèle optimisé est intégré
        au code embarqué sous forme de tableau d'octets statique stocké dans la
        mémoire flash.
  \item \textbf{Exécution de l'inférence} : le microcontrôleur exécute le modèle
        sur les données acquises par ses capteurs, sans connexion réseau. Les
        résultats déclenchent des actions locales.
\end{enumerate}

\textbf{Caractéristique fondamentale :} le modèle est figé au moment de la
compilation. Le microcontrôleur n'effectue aucune forme d'apprentissage ; il
applique les calculs définis.

% ---------------------------------------------------------------------
\section{Techniques de compression et d'optimisation des modèles}

\subsection{Introduction à la compression}

La compression des modèles représente un domaine essentiel pour le déploiement
de l'IA sur dispositifs embarqués~\cite{sze2020efficient}. L'objectif consiste
à réduire significativement la taille des modèles tout en minimisant la
dégradation des performances. Cette réduction permet de respecter les
contraintes mémorielles, d'améliorer l'efficacité énergétique et la vitesse
d'exécution.

Les techniques peuvent être classées en plusieurs catégories : quantification,
réduction de complexité structurelle (élagage, distillation), factorisation et
approches hybrides~\cite{han2016deep}.

\subsection{La quantification (\emph{quantization})}

\subsubsection{Définition et principes}

La quantification est la technique de compression la plus répandue et la plus
efficace en TinyML~\cite{jacob2018quantization,gholami2021survey}. Elle réduit
la précision des représentations numériques des poids et des activations du
réseau.

\textbf{Contexte théorique :} en apprentissage classique, les réseaux
manipulent des valeurs en virgule flottante 32~bits (Float32), offrant une
grande précision mais au prix d'un espace mémoire important (4~octets par
nombre). Dans l'embarqué, cette précision est généralement superflue.

\textbf{Principe :} la quantification consiste à représenter les valeurs sur
8~bits entiers (Int8), divisant par quatre la taille mémoire requise. Un modèle
de 4~Mo en Float32 n'occupera plus que 1~Mo en Int8.

\subsubsection{Mécanismes}

La quantification implique un processus de mise en correspondance
(\emph{mapping}) préservant l'information. Avant d'expliciter ce mécanisme, il
est nécessaire de définir ce qu'est un \textbf{tenseur}.

\paragraph{Qu'est-ce qu'un tenseur ?} En apprentissage automatique, un tenseur
est une généralisation des matrices et des vecteurs à un nombre quelconque de
dimensions. C'est la structure de données fondamentale utilisée pour
représenter les poids, les activations et les données d'entrée dans les réseaux
de neurones. Un tenseur d'ordre $n$ est un tableau à $n$ indices ; chaque
élément est identifié par $n$ coordonnées $(i_1, i_2, \ldots, i_n)$.

\begin{itemize}
  \item \textbf{Ordre 0} (scalaire) : une simple valeur numérique, par exemple $3{,}14$ ;
  \item \textbf{Ordre 1} (vecteur) : un tableau 1D, par exemple les biais d'une couche : $[0{,}2;\ -0{,}5;\ 1{,}3;\ 0{,}8]$ ;
  \item \textbf{Ordre 2} (matrice) : un tableau 2D, comme les poids d'une couche entièrement connectée ;
  \item \textbf{Ordre 3} : un cube de données, typiquement les activations d'une couche convolutionnelle avec plusieurs canaux ;
  \item \textbf{Ordre 4} : les poids des couches convolutionnelles (hauteur $\times$ largeur $\times$ canaux d'entrée $\times$ canaux de sortie).
\end{itemize}

Dans le contexte du TinyML, les \textbf{poids} du réseau sont stockés sous forme
de tenseurs (généralement d'ordre 2 ou 4), les \textbf{activations} (sorties
intermédiaires de chaque couche) sont également des tenseurs, et la
\textbf{quantification} s'applique à chaque tenseur individuellement, avec ses
propres paramètres (\emph{scale} et \emph{zero\_point}).

\paragraph{Formulation mathématique.} La relation fondamentale de la
quantification affine s'exprime par~\cite{jacob2018quantization} :
\begin{equation}
  r \approx S \times (q - Z)
  \label{eq:quant}
\end{equation}
où $r$ est la valeur réelle (Float32), $q$ la valeur quantifiée (entier Int8
dans $[-128, 127]$), $S$ le facteur d'échelle (\emph{scale}, réel positif) et
$Z$ le point zéro (\emph{zero\_point}, entier de même type que $q$).

\paragraph{Processus détaillé.}
\begin{enumerate}
  \item \textbf{Analyse de la plage dynamique} : pour chaque tenseur, on
        détermine les valeurs minimale et maximale ;
  \item \textbf{Calcul des paramètres} : on calcule $S$ et $Z$ pour faire
        correspondre la plage du tenseur à $[-128, 127]$ ;
  \item \textbf{Application de la quantification} : chaque valeur est convertie
        selon la formule~\eqref{eq:quant} inversée ;
  \item \textbf{Exécution en arithmétique entière} : pendant l'inférence, les
        multiplications se font en Int8 et les accumulations en Int32.
\end{enumerate}

\paragraph{Exemple détaillé.} Considérons un tenseur de poids d'une couche
dense avec les valeurs suivantes (en Float32) :
\[
  W = [0{,}23;\ -0{,}45;\ 0{,}12;\ -0{,}87;\ 0{,}65;\ -0{,}33]
\]

\textit{Étape 1 -- Détermination de la plage :} $\min = -0{,}87$,
$\max = 0{,}65$, d'où $\text{plage} = 0{,}65 - (-0{,}87) = 1{,}52$.

\textit{Étape 2 -- Calcul du facteur d'échelle :} sur 8~bits, l'intervalle
$[-128, 127]$ compte 255 pas :
\[
  S = \frac{\text{plage}}{255} = \frac{1{,}52}{255} \approx 0{,}0059608
\]

\textit{Étape 3 -- Calcul du point zéro :} on choisit $Z$ de sorte que la
valeur minimale corresponde à $-128$ :
\[
  Z = -128 - \operatorname{round}\!\left(\frac{\min}{S}\right)
    = -128 - \operatorname{round}(-145{,}95) = -128 + 146 = 18
\]

\textit{Étape 4 -- Quantification de chaque valeur :}
\[
  q_i = \operatorname{round}\!\left(\frac{w_i}{S}\right) + Z
  \quad\Rightarrow\quad
  W_q = [57;\ -57;\ 38;\ -128;\ 127;\ -37]
\]
Toutes les valeurs sont bien comprises dans l'intervalle Int8.

\textit{Étape 5 -- Vérification par déquantification :} pour la première valeur,
\[
  r \approx 0{,}0059608 \times (57 - 18) = 0{,}0059608 \times 39 \approx 0{,}2325
\]
soit une erreur de $0{,}0025$ par rapport à la valeur d'origine $0{,}23$.

\begin{table}[H]
\centering
\small
\begin{tabular}{rrrr}
\toprule
Valeur d'origine & Valeur quantifiée & Valeur déquantifiée & Erreur absolue \\
\midrule
$0{,}23$  & $57$   & $0{,}2325$  & $0{,}0025$ \\
$-0{,}45$ & $-57$  & $-0{,}4471$ & $0{,}0029$ \\
$0{,}12$  & $38$   & $0{,}1192$  & $0{,}0008$ \\
$-0{,}87$ & $-128$ & $-0{,}8703$ & $0{,}0003$ \\
$0{,}65$  & $127$  & $0{,}6497$  & $0{,}0003$ \\
$-0{,}33$ & $-37$  & $-0{,}3278$ & $0{,}0022$ \\
\bottomrule
\end{tabular}
\caption{Exemple de quantification Int8 d'un tenseur de six poids ($S \approx 0{,}0059608$, $Z = 18$).}
\label{tab:quant-exemple}
\end{table}

\textbf{Conclusion de l'exemple :} l'erreur d'arrondi est toujours inférieure
à $S/2 \approx 0{,}003$. Elle est généralement acceptable pour les réseaux de
neurones : la perte de précision globale du modèle reste le plus souvent de
l'ordre de 1 à 2~points~\cite{jacob2018quantization,krishnamoorthi2018}.

\subsubsection{Avantages}
\begin{itemize}
  \item \textbf{Réduction de taille} : division par quatre de la taille des
        poids. Pour un modèle de 20~Mo, cela représente une économie de 15~Mo.
  \item \textbf{Amélioration des performances} : de nombreux microcontrôleurs
        exécutent les opérations entières bien plus rapidement que les opérations
        flottantes. Les opérations Int8 sont particulièrement efficaces sur les
        architectures disposant d'instructions SIMD.
  \item \textbf{Réduction de la consommation énergétique} : une addition 8~bits
        entière consomme environ 30 fois moins d'énergie qu'une addition
        flottante 32~bits, et une multiplication environ 18 fois
        moins~\cite{horowitz2014energy}.
  \item \textbf{Compatibilité avec le matériel} : de nombreux accélérateurs d'IA
        (NPU) sont optimisés pour les calculs Int8.
\end{itemize}

\subsubsection{Méthodes de quantification}

\textbf{Quantification post-entraînement (PTQ)} : convertit un modèle déjà
entraîné. Rapide et simple, elle suffit dans la majorité des cas. La perte de
précision est généralement inférieure à 1--2~\% pour la
classification~\cite{krishnamoorthi2018}.

\textbf{Quantification consciente de l'entraînement (QAT)} : intègre la
quantification dès la phase d'entraînement via une \emph{fake quantization}.
Elle permet d'atteindre des quantifications plus agressives (Int4, binaire)
avec une meilleure préservation de la précision~\cite{jacob2018quantization}.

\begin{lstlisting}[language=Python, caption={Quantification post-entraînement Int8 avec TensorFlow Lite.}, label={lst:ptq}]
import tensorflow as tf

converter = tf.lite.TFLiteConverter.from_saved_model(saved_model_dir)
converter.optimizations = [tf.lite.Optimize.DEFAULT]
# Jeu de données représentatif pour calibrer scale et zero_point
converter.representative_dataset = generateur_donnees_representatives
converter.target_spec.supported_ops = [tf.lite.OpsSet.TFLITE_BUILTINS_INT8]
converter.inference_input_type = tf.int8
converter.inference_output_type = tf.int8
quantized_model = converter.convert()
open("modele_int8.tflite", "wb").write(quantized_model)
\end{lstlisting}

\subsubsection{Limites et défis}
\begin{itemize}
  \item \textbf{Perte de précision} : généralement faible pour la
        classification, elle peut être significative pour la détection d'objets
        ou les opérations sensibles.
  \item \textbf{Gestion des valeurs aberrantes} (\emph{outliers}) : les valeurs
        extrêmes élargissent la plage et dégradent la résolution des autres.
  \item \textbf{Solutions} : quantification par canal, écrêtage
        (\emph{clipping}), utilisation de la QAT.
  \item \textbf{Compatibilité} : certains opérateurs ne sont pas disponibles en
        Int8 dans les moteurs d'inférence embarqués.
\end{itemize}

\subsubsection{Ordres de grandeur}

\begin{table}[H]
\centering
\small
\begin{tabular}{lcccc}
\toprule
Modèle & Taille Float32 & Taille Int8 & Perte de précision & Gain de vitesse \\
\midrule
MobileNetV1  & 16,9 Mo & 4,3 Mo  & $-0{,}8$ \% & $\times 3{,}2$ \\
MobileNetV2  & 14,0 Mo & 3,5 Mo  & $-1{,}2$ \% & $\times 4{,}1$ \\
InceptionV3  & 92,0 Mo & 23,0 Mo & $-1{,}5$ \% & $\times 2{,}8$ \\
\bottomrule
\end{tabular}
\caption{Comparaison indicative des performances avant et après quantification Int8 (ordres de grandeur ; voir~\cite{krishnamoorthi2018,jacob2018quantization}).}
\label{tab:bench-quant}
\end{table}

\subsection{L'élagage (\emph{pruning})}

\subsubsection{Définition et motivation}

L'élagage est une technique de compression structurelle qui supprime les
connexions non essentielles~\cite{han2016deep}. Les réseaux entraînés
contiennent de nombreuses connexions dont les poids sont proches de zéro et qui
contribuent peu à la décision finale.

\textbf{Origine biologique :} l'élagage s'inspire des mécanismes de plasticité
synaptique, dans lesquels les connexions les moins utilisées sont éliminées.

\subsubsection{Méthodologies}

\textbf{Élagage non structuré :} suppression des poids individuellement en
fonction de leur magnitude. Le taux de compression peut être très élevé
(90--95~\%), mais il nécessite des formats de stockage creux spécifiques pour
se traduire par un gain réel.

\textbf{Élagage structuré :} suppression de blocs entiers (filtres, neurones,
canaux). Il réduit effectivement le nombre d'opérations et reste compatible
avec le matériel existant.

\textbf{Algorithme typique :} (1) entraînement initial ; (2) tri des poids par
magnitude ; (3) sélection des $N$~\% de poids les plus petits ; (4) mise à
zéro ; (5) réentraînement court (\emph{fine-tuning}).

\begin{lstlisting}[language=Python, caption={Élagage fondé sur un critère d'importance (approximation de Taylor).}, label={lst:pruning}]
import numpy as np

def taylor_pruning(weights, gradients, pruning_rate=0.5):
    importance = np.abs(weights * gradients)
    threshold = np.percentile(importance, pruning_rate * 100)
    mask = importance > threshold
    return weights * mask, mask
\end{lstlisting}

\subsubsection{Processus complet}

Le processus enchaîne pré-entraînement, élagage (structuré ou non),
\emph{fine-tuning} et itération progressive. \textbf{Stratégie
recommandée :} plusieurs itérations avec des taux modérés (par exemple 20~\% à
chaque itération) plutôt qu'une seule itération avec un taux élevé.

\subsubsection{Résultats expérimentaux}

\begin{table}[H]
\centering
\small
\begin{tabular}{cccc}
\toprule
Taux d'élagage & Précision & Perte & Réduction des paramètres \\
\midrule
20 \% & 70,9 \% & $-0{,}3$ \% & 20 \% \\
40 \% & 70,1 \% & $-1{,}1$ \% & 40 \% \\
60 \% & 68,5 \% & $-2{,}7$ \% & 60 \% \\
80 \% & 65,2 \% & $-6{,}0$ \% & 80 \% \\
\bottomrule
\end{tabular}
\caption{Impact indicatif de l'élagage sur la précision de MobileNetV1 (ImageNet).}
\label{tab:pruning}
\end{table}

\subsection{La distillation (\emph{knowledge distillation})}

\subsubsection{Concept}

La distillation transfère les connaissances d'un modèle complexe (le
\emph{professeur}) vers un modèle plus petit (l'\emph{élève}), permettant de
créer des modèles compacts tout en conservant une grande partie des
performances~\cite{hinton2015distilling}.

\textbf{Analogie pédagogique :} un expert transmet non seulement les réponses
finales, mais également son raisonnement à un élève.

\subsubsection{Principe théorique}

Le transfert de connaissances peut porter sur les sorties (\emph{output
knowledge}), sur les représentations intermédiaires (\emph{intermediate
knowledge}) ou sur les relations entre exemples (\emph{relation knowledge}).
La fonction de perte combine deux termes :
\begin{equation}
  L_{\text{total}} = \alpha \, L_{\text{classification}} + \beta \, L_{\text{distillation}}
\end{equation}
Une \textbf{température} $T > 1$ est utilisée pour adoucir les distributions de
probabilités produites par le professeur~\cite{hinton2015distilling} :
\begin{equation}
  P_i = \frac{\exp(z_i / T)}{\sum_j \exp(z_j / T)}
\end{equation}

\subsubsection{Méthodes}
\begin{itemize}
  \item \textbf{Distillation standard} : l'élève apprend à reproduire les sorties du professeur.
  \item \textbf{Distillation par représentations intermédiaires} : l'élève apprend les représentations de plusieurs couches internes.
  \item \textbf{Distillation mutuelle} : plusieurs élèves s'enseignent mutuellement.
  \item \textbf{Auto-distillation} : un modèle enseigne à une version plus petite de lui-même.
\end{itemize}

\subsubsection{Résultats}

\begin{table}[H]
\centering
\small
\begin{tabular}{llcccc}
\toprule
Professeur & Élève & Param. P & Param. E & Précision P & Précision E \\
\midrule
ResNet-152 & ResNet-18 & 60,2 M & 11,7 M & 78,3 \% & 74,8 \% \\
ResNet-152 & MobileNet & 60,2 M & 4,2 M  & 78,3 \% & 71,3 \% \\
\bottomrule
\end{tabular}
\caption{Comparaison indicative des performances (top-1 ImageNet) avant et après distillation ; M = millions de paramètres.}
\label{tab:distillation}
\end{table}

\subsection{Combinaison des techniques}

Les techniques sont complémentaires et sont généralement appliquées dans
l'ordre suivant~\cite{han2016deep} : (1) distillation pour réduire la taille de
l'architecture ; (2) élagage pour réduire la complexité ; (3) quantification
pour réduire la précision numérique ; avec un \emph{fine-tuning} après chaque
étape.

\begin{figure}[H]
\centering
\begin{tikzpicture}[node distance=0.5cm]
  \node[bloco, minimum width=2.3cm] (m) {Modèle Float32\\\footnotesize(plusieurs Mo)};
  \node[bloc, minimum width=2.1cm, right=of m] (d) {Distillation};
  \node[bloc, minimum width=2.1cm, right=of d] (e) {Élagage\\+ \emph{fine-tuning}};
  \node[bloc, minimum width=2.1cm, right=of e] (q) {Quantification\\Int8};
  \node[blocv, minimum width=2.3cm, right=of q] (f) {Modèle TinyML\\\footnotesize(quelques dizaines de Ko)};
  \draw[fleche] (m) -- (d);
  \draw[fleche] (d) -- (e);
  \draw[fleche] (e) -- (q);
  \draw[fleche] (q) -- (f);
\end{tikzpicture}
\caption{Enchaînement typique des techniques de compression.}
\label{fig:compression}
\end{figure}

\begin{table}[H]
\centering
\small
\begin{tabular}{lccc}
\toprule
Méthode & Réduction de taille & Perte de précision & Gain de vitesse \\
\midrule
Quantification seule & $\times 4$         & $< 1$ \%     & $\times 2$ à $4$ \\
Élagage seul         & $\times 5$ à $10$  & 2 à 5 \%     & $\times 1$ à $2$ \\
Distillation seule   & $\times 10$ à $50$ & 3 à 8 \%     & $\times 1$ à $3$ \\
Combinaison          & $\times 100$ à $300$ & 5 à 15 \%  & $\times 5$ à $15$ \\
\bottomrule
\end{tabular}
\caption{Efficacité comparée des méthodes de compression (ordres de grandeur).}
\label{tab:compression}
\end{table}

\subsection{Techniques avancées}
\begin{itemize}
  \item \textbf{Factorisation de matrices} : décomposition des matrices de poids
        en produits de matrices de plus petites dimensions (SVD, rang faible).
  \item \textbf{Binarisation} : représentation des poids et activations par des
        valeurs binaires ($-1$ ou $+1$). La réduction de taille est extrême
        ($\times 32$) mais la perte de précision importante.
  \item \textbf{Réseaux creux} : utilisation de représentations creuses pour les
        matrices de poids (formats COO, CSR).
  \item \textbf{Recherche automatique d'architecture (NAS)} : conception
        automatique de réseaux adaptés à l'embarqué, comme MobileNet,
        EfficientNet ou MCUNet~\cite{lin2020mcunet}.
\end{itemize}

\subsection{Métriques d'évaluation}
\begin{itemize}
  \item \textbf{Taux de compression} : $\text{Taille}_{\text{originale}} / \text{Taille}_{\text{compressée}}$.
  \item \textbf{Métriques de performance} : précision, latence, consommation énergétique, mémoire RAM et flash.
  \item \textbf{Bancs d'essai} : ImageNet, CIFAR-10/100 et, pour l'embarqué, MLPerf Tiny~\cite{banbury2021mlperftiny}.
\end{itemize}

\subsection{Études de cas}
\begin{itemize}
  \item \textbf{Détection de mots-clés} : un modèle CNN de l'ordre du mégaoctet
        peut être ramené à une centaine de kilooctets, avec une précision
        supérieure à 90~\% et une latence de quelques dizaines de millisecondes
        sur microcontrôleur~\cite{banbury2021mlperftiny}.
  \item \textbf{Classification d'images} : des réseaux de type MobileNet ou
        MCUNet, quantifiés en Int8, tiennent dans moins de 1~Mo et s'exécutent
        sur des microcontrôleurs haut de gamme~\cite{lin2020mcunet}.
  \item \textbf{Séries temporelles} (vibrations, capteurs inertiels) : de petits
        réseaux denses ou auto-encodeurs de quelques dizaines de kilooctets
        suffisent pour la détection d'anomalies, avec une latence de quelques
        millisecondes sur ESP32 ; c'est le type de modèle visé par le scénario illustratif de la partie~\ref{ch:fondements}.
\end{itemize}
